\documentclass[11pt]{article}
\usepackage[margin=1in]{geometry}
\usepackage{amsmath,amssymb,amsthm,hyperref}
\newtheorem{theorem}{Theorem}
\newtheorem{lemma}{Lemma}
\newcommand{\F}{\mathbb F_2}
\title{A cube counterexample to the complementation-set correspondence}
\author{Principia Math In Progress}
\date{29 September 2026}
\begin{document}
\maketitle

\paragraph{Task 1 (original problem).}
MathDB \href{https://mathdb.com/p/380333}{380333} records
Genest's Conjecture 1.7.17 \cite[p.~41]{Genest}: for a bicolored graph
without isolated vertices or twins, its parity class is in bijection
with its complementation sets. We disprove the conjecture, including
for a graph with its natural coloring.

\begin{theorem}\label{main}
Let \(G\) be the three-dimensional cube, all of whose vertices are black.
Then \(G\) is connected, has no twins or isolated vertices, has exactly
\(56\) complementation sets, and has exactly \(28\) graphs in its parity
class. Consequently no bijection in the conjecture exists.
\end{theorem}

\paragraph{Definitions from the source.}
Local complementation at \(u\) toggles the edges between distinct
neighbors of \(u\). If \(u\) is white, it also reverses its neighbors'
colors; if \(u\) is black, all colors stay unchanged.
Allowed parity steps are a white singleton \(u\), or
\([uv]=uvu\) for adjacent black vertices \(u,v\).
A parity word is any finite sequence of allowed steps, including the
empty word. The parity class \([G]\) consists of the resulting colored
graphs on the original vertex set.
A reduced word uses disjoint singleton or pair supports; its support
is a \emph{complementation set}. The result for support \(S\) is denoted
\(G^S\) and is independent of the reduced word.
These are Definitions 1.6.1--1.6.4 and Theorem 1.6.3 in \cite{Genest}.
Twins mean distinct vertices with equal open neighborhoods or equal
closed neighborhoods. The natural coloring makes odd degrees black
and even degrees white.

Starting with an all-black graph, every parity step is an adjacent
black pair, and every resulting graph is all black. We work entirely
in this case. All matrices below are over \(\F\).

\begin{lemma}\label{pivots}
Let \(A\) be the adjacency matrix of an all-black graph.
Its complementation sets are exactly the sets \(S\) for which the
principal submatrix \(A[S]\) is nonsingular, with the empty matrix
regarded as nonsingular. Every graph in its parity class equals
\(G^S\) for such a set.
\end{lemma}
\begin{proof}
For a nonsingular principal block \(P=A[S]\), write
\[
 A=\begin{pmatrix}P&Q\\Q^t&R\end{pmatrix}.
\]
Solving the first block of \(y=Ax\) for \(x_S\) gives
\begin{equation}\label{ppt}
 \begin{pmatrix}x_S\\y_{S^c}\end{pmatrix}
 =
 \underbrace{\begin{pmatrix}
 P^{-1}&P^{-1}Q\\Q^tP^{-1}&R+Q^tP^{-1}Q
 \end{pmatrix}}_{A*S}
 \begin{pmatrix}y_S\\x_{S^c}\end{pmatrix}.
\end{equation}
This defines the principal pivot \(A*S\).

For an adjacent black pair \(u,v\), its block is
\(P=\left(\begin{smallmatrix}0&1\\1&0\end{smallmatrix}\right)=P^{-1}\).
If its outside incidence rows are \(c^t,d^t\), formula \eqref{ppt}
swaps these rows and adds \(cd^t+dc^t\) to the outside block.
The source's three local complementations \(uvu\) do exactly this:
outside adjacencies \((c_i,d_i)\) become \((d_i,c_i)\), and the
toggle between outside vertices \(i,j\) is
\[
 c_ic_j+(c_i+d_i)(c_j+d_j)+d_id_j=c_id_j+d_ic_j.
\]
All colors stay black.

For completeness, let
\[
 L_A=\{(x,Ax):x\in\F^n\},
\]
and let \(\sigma_S\) interchange \(x_i\) and \(y_i\) for \(i\in S\).
These swaps commute and satisfy
\(\sigma_S\sigma_T=\sigma_{S\mathbin{\triangle}T}\).
Writing \(D_S\) for the diagonal indicator of \(S\), the first
coordinate of \(\sigma_S(x,Ax)\) is
\[
 C_Sx,\qquad C_S=I+D_S+D_SA.
\]
Its rows outside \(S\) are identity rows and its rows inside \(S\)
are rows of \(A\). Therefore
\(\det C_S=\det A[S]\).
It follows that \(\sigma_SL_A\) is the graph of a matrix precisely
when \(A[S]\) is nonsingular; that matrix is \(A*S\).

Any parity word is a sequence of pair swaps. Its final subspace is
\(\sigma_SL_A\), where \(S\) is the symmetric difference of all its
pair supports. The final graph is a matrix graph, so \(A[S]\) is
nonsingular and the final matrix is \(A*S\).

Conversely, if \(A[S]\) is nonsingular and nonempty, it is alternating:
it is symmetric with zero diagonal. It has a nonzero off-diagonal
entry, hence an allowed black pair. Its Schur complement on the
unused part of \(S\) is again alternating and nonsingular, by
the displayed formula and determinant factorization. Repeating gives
disjoint pair steps covering \(S\). This is a reduced parity word,
and \eqref{ppt} shows that it realizes \(A*S\).
\end{proof}

\begin{lemma}\label{fibers}
Suppose that \(A\) is the adjacency matrix of a connected bipartite
graph and \(A^2=I\). For any two complementation sets \(S,T\),
\[
 G^S=G^T\quad\Longleftrightarrow\quad
 T=S\ \hbox{or}\ T=V(G)\setminus S.
\]
In particular, every graph in the parity class has exactly two
complementation sets.
\end{lemma}
\begin{proof}
By Lemma~\ref{pivots}, equality is equivalent to
\(\sigma_{S\mathbin{\triangle}T}L_A=L_A\).
Let \(D\) be the diagonal indicator of a set \(U\) with
\(\sigma_UL_A=L_A\). The two coordinates of the swapped vector are
\[
 x'=(I+D+DA)x,\qquad y'=(D+(I+D)A)x.
\]
Because \(y'=Ax'\) for every \(x\), we obtain
\begin{equation}\label{stabilizer}
 D+DA+AD+ADA=0.
\end{equation}
For adjacent vertices \(i,j\), bipartiteness implies that they have
no common neighbor, so \((ADA)_{ij}=0\). The \((i,j)\) entry of
\eqref{stabilizer} is therefore \(d_i+d_j=0\).
Connectivity forces all the \(d_i\) to be equal.
Thus only \(U=\varnothing\) or \(U=V(G)\) can fix \(L_A\).
Both do fix it: the empty swap is the identity, and the full swap
gives \(L_{A^{-1}}=L_A\) because \(A^2=I\).

Consequently equality of two pivot results holds exactly when
\(S\mathbin{\triangle}T\) is empty or full.
Moreover the complement of every feasible \(S\) is feasible, since
\(\sigma_{V\setminus S}L_A=\sigma_SL_A\) is a matrix graph.
The two supports are distinct, completing the proof.
\end{proof}

\begin{proof}[Proof of Theorem~\ref{main}]
Label the bipartition as
\[
 L=\{\ell_1,\ell_2,\ell_3,\ell_4\},\qquad
 R=\{r_1,r_2,r_3,r_4\},
\]
with \(\ell_i\) adjacent to \(r_j\) precisely when \(i\ne j\).
This is the cube: the labels can be taken as
\[
 (\ell_1,\ell_2,\ell_3,\ell_4)=(000,011,101,110),\qquad
 (r_1,r_2,r_3,r_4)=(111,100,010,001).
\]
Edges then join binary strings differing in one coordinate.
The graph is connected and every vertex has degree three.
Thus the all-black coloring is its natural coloring.

Two vertices in the same part omit different opposite vertices from
their open neighborhoods, so these neighborhoods differ. Open
neighborhoods of vertices in different parts are nonempty and lie
in disjoint parts, so cannot agree either. Equal closed neighborhoods
of distinct vertices would require them to be adjacent. But for
an adjacent pair, one closed neighborhood contains one left vertex
and the other contains three left vertices. Hence there are no twins.

Let \(J_4\) be the all-ones matrix and put \(B=I_4+J_4\).
Then
\[
 A=\begin{pmatrix}0&B\\B&0\end{pmatrix},\qquad
 J_4^2=0,\qquad B^2=I_4,\qquad A^2=I_8.
\]
It remains to count the nonsingular principal submatrices.
For \(S=X\cup Y\), with \(X\subseteq L\), \(Y\subseteq R\),
the principal matrix is
\(\left(\begin{smallmatrix}0&B[X,Y]\\B[X,Y]^t&0\end{smallmatrix}\right)\).
It can be nonsingular only if \(|X|=|Y|=k\), and in that case
it is nonsingular exactly when \(B[X,Y]\) is.
Identify \(X,Y\) with their index sets in \(\{1,2,3,4\}\).
The possibilities are:
\[
\begin{array}{c|c|c}
 k&\text{condition for nonsingularity}&\text{number}\\ \hline
 0&\text{empty}&1\\
 1&X\ne Y&12\\
 2&|X\cap Y|=1\text{ or }2&24+6=30\\
 3&X\ne Y&12\\
 4&X=Y=\{1,2,3,4\}&1
\end{array}
\]
For \(k=2\), disjoint sets give an all-ones matrix of rank one;
intersection sizes one and two give, up to row and column
permutations, respectively
\[
 \begin{pmatrix}0&1\\1&1\end{pmatrix},\qquad
 \begin{pmatrix}0&1\\1&0\end{pmatrix},
\]
both of determinant one.
For \(k=3\), equal sets give \(I_3+J_3\), which kills the all-ones
vector, while unequal sets give, up to permutations,
\[
 \begin{pmatrix}0&1&1\\1&0&1\\1&1&1\end{pmatrix}
\]
of determinant one. The other rows are immediate, using \(B^2=I_4\)
for the last row.
There are therefore \(1+12+30+12+1=56\) complementation sets.
Lemma~\ref{fibers} gives precisely \(56/2=28\) parity-class graphs.
This strict cardinality discrepancy excludes every possible bijection,
not only the canonical map \(S\mapsto G^S\).
\end{proof}

\paragraph{An infinite counterfamily.}
Every odd-dimensional cube \(Q_d\), \(d\ge3\), with its natural
all-black coloring has the same two-to-one correspondence.
Indeed, write \(A=P_1+\cdots+P_d\), where \(P_i\) flips coordinate
\(i\). The commuting involutions \(P_i\) satisfy
\(A^2=dI=I\) over \(\F\). The graph is connected and bipartite.
Distinct vertices have at most two common neighbors, whereas every
vertex has \(d\ge3\) neighbors, so open twins are impossible.
Closed twins are impossible too: they would be adjacent, but an
adjacent pair has no common neighbor and degree at least three.
Lemma~\ref{fibers} applies.

\paragraph{Verification and scope.}
The exact verifier separately explores unrestricted parity words and
reduced words using the source's three local complementations.
It checks all \(256\) supports against independent binary elimination,
finds \(56\) supports and \(28\) distinct labeled colored graphs,
and confirms that every fiber is a complementary pair.
For the labels above, one full reduced word returning to the original
graph is
\[
 [\ell_1r_2][\ell_2r_1][\ell_3r_4][\ell_4r_3].
\]
The proof above is independent of this computation.
The conjecture fails under its full source definitions, with no
transcription repair or missing-hypothesis counterexample involved.

\paragraph{Task 2 (proposed follow-up).}
Determine the number of complementation sets of \(Q_d\) for odd
\(d\ge5\), and hence the cardinality of its parity class.

\paragraph{Task 3 (proposed follow-up).}
Determine the smallest order of a naturally colored graph without
isolated vertices or twins for which the correspondence is not injective.

\paragraph{Task 4 (proposed follow-up).}
Classify the connected bipartite graphs whose binary adjacency matrix
satisfies \(A^2=I\), by their degree and common-neighbor constraints.

\paragraph{Task 5 (proposed follow-up).}
For general bicolored graphs, classify the diagonal indicator matrices
\(D\) satisfying \(D+DM+MD+MDM=0\), where \(M\) has diagonal one
at white vertices, and determine the possible dimensions of this
space for graphs without twins.

\begin{thebibliography}{1}
\bibitem{Genest}
F. Genest, \emph{Graphes eul\'eriens et compl\'ementarit\'e locale}
(English arXiv title: \emph{Eulerian graphs and local complementation}),
Ph.D. thesis, Universit\'e de Montr\'eal;
arXiv:math/0701421v1 (2007), 103 pages.
Definitions 1.6.1--1.6.4, Theorem 1.6.3,
and Conjecture 1.7.17.
\url{https://arxiv.org/abs/math/0701421}.
\end{thebibliography}
\end{document}
